\documentclass[a4paper,11pt]{article}
\usepackage{fontspec}
\IfFontExistsTF{HaranoAjiMincho-Regular}{%
  \setmainfont{HaranoAjiMincho-Regular}[
    BoldFont=HaranoAjiMincho-Bold,ItalicFont=HaranoAjiMincho-Regular,
    BoldItalicFont=HaranoAjiMincho-Bold]
  \setsansfont{HaranoAjiGothic-Medium}[
    BoldFont=HaranoAjiGothic-Bold,ItalicFont=HaranoAjiGothic-Medium]
}{%
  \setmainfont{Noto Serif CJK JP}[
    ItalicFont={Noto Serif CJK JP},BoldItalicFont={Noto Serif CJK JP Bold}]
  \setsansfont{Noto Sans CJK JP}[ItalicFont={Noto Sans CJK JP}]
}
\setmonofont{Latin Modern Mono}
\XeTeXlinebreaklocale "ja"
\XeTeXlinebreakskip=0pt plus 1pt
\usepackage{amsmath,amssymb,amsthm,mathtools}
\usepackage[margin=25mm]{geometry}
\usepackage{enumitem,array,booktabs}
\usepackage[unicode,colorlinks=true,linkcolor=blue,citecolor=blue,urlcolor=blue]{hyperref}
\usepackage{bookmark}
\allowdisplaybreaks[2]
\setlength{\emergencystretch}{2em}
\numberwithin{equation}{section}
\newtheorem{theorem}{定理}[section]
\newtheorem{proposition}[theorem]{命題}
\newtheorem{lemma}[theorem]{補題}
\newtheorem{corollary}[theorem]{系}
\theoremstyle{definition}
\newtheorem{definition}[theorem]{定義}
\newtheorem{example}[theorem]{例}
\newtheorem{remark}[theorem]{注意}
\renewcommand{\proofname}{証明}
\renewcommand{\contentsname}{目次}
\renewcommand{\refname}{参考文献}
\newcommand{\C}{\mathbb C}
\DeclareMathOperator{\Ext}{Ext}
\DeclareMathOperator{\Tor}{Tor}
\DeclareMathOperator{\gldim}{gldim}
\DeclareMathOperator{\GKdim}{GKdim}
\DeclareMathOperator{\ASind}{ASindex}
\DeclareMathOperator{\Aut}{Aut}
\DeclareMathOperator{\id}{id}
\DeclareMathOperator{\im}{im}
\DeclareMathOperator{\pd}{pd}
\IfFileExists{build/build_info.tex}{\input{build/build_info.tex}}{%
  \newcommand{\BuildTimestamp}{\texttt{python3 scripts/build.py} により生成}
  \newcommand{\BuildFilename}{\texttt{python3 scripts/build.py} により生成}}
\hypersetup{pdftitle={最小生成次数 (1,1,1,2) の AS-regular 代数と Ore 拡大}}
\begin{document}
\begin{titlepage}
\centering
\vspace*{30mm}
{\LARGE 最小生成次数 \((1,1,1,2)\) の\\[3mm]
AS-regular 代数と Ore 拡大\par}
\vspace{10mm}
{\large 中心拡大と重み \((1^n,2)\) への一般化\par}
\vspace{20mm}
\begin{minipage}{0.85\textwidth}
連結次数付き両側 noetherian AS-regular 代数を扱う。
4 次元で最小生成次数が \((1,1,1,2)\) の場合を、
3 次元 quadratic AS-regular 代数上の次数 \(2\) の Ore 拡大として記述する。
3 次元 \((1,1,2)\) 型の regular central extension を特徴付け、
一般の \((1^n,2)\) 型については AS index の仮定の下で同じ構造定理を証明する。
最小生成次数だけでは結論が成り立たない例も与える。
\end{minipage}
\vfill
{\small 作成日時：\BuildTimestamp\par}
\vspace{3mm}
{\small PDF ファイル名：\texttt{\BuildFilename}\par}
\vspace{8mm}
{\small \url{https://github.com/uedakazushi/nc-p1112}\par}
\end{titlepage}
\tableofcontents
\clearpage
\section{対象と結果の範囲}\label{sec:conventions}

係数体を \(k=\C\) とする。代数は連結局所有限な
\(\mathbb N\)-次数付き \(k\)-代数
\(A=\bigoplus_{j\ge0}A_j\)、\(A_0=k\) とする。
\(A_+=\bigoplus_{j>0}A_j\)、\(k=A/A_+\) と書き、
shift は \(M(a)_j=M_{j+a}\) とする。
\(H_M(t)=\sum_j(\dim_kM_j)t^j\) は、各次数有限次元で次数下有界な
次数付きベクトル空間 \(M\) の Hilbert 級数である。

AS-regular とは、大域次元 \(d\) と GK 次元が有限であり、
左右の AS-Gorenstein 条件を満たすことをいう。
右加群の規約では、その条件は
\[
\underline{\Ext}^{i}_A(k,A)\cong
\begin{cases}k(\ell)&i=d,\\0&i\ne d\end{cases}
\]
である。\(\underline{\Ext}\) はすべての内部次数を含む次数付き Ext を表す。
\(\ell=\ASind A\) を AS index と呼ぶ。
本文の構造定理では、両側 noetherian 性を仮定する。

\emph{最小生成次数が \((1^n,2)\)} とは、
\(A_+/A_+^2\) が次数 \(1\) に \(n\) 次元、次数 \(2\) に \(1\) 次元を持ち、
他の次数では零であることをいう。
特に、次数 \(2\) の生成元は次数 \(1\) の元の積で消去できない。
\[
V=A_1,\qquad B=k\langle A_1\rangle\subseteq A
\]
と書く。\(B\) はこの定義から \(A\) に内在的に定まる。
自由テンソル代数を \(T(V)=\bigoplus_{j\ge0}V^{\otimes_k j}\)、
その次数 \(j\) の部分を \(T_j(V)=V^{\otimes_k j}\) と書く。

次数 \(2\) の Ore 拡大の規約は
\[
A=B[w;\sigma,\delta],\qquad
wb=\sigma(b)w+\delta(b),\qquad |w|=2,
\]
である。ここで \(\sigma\in\Aut_{\mathrm{gr}}B\)、
\(\delta(B_j)\subseteq B_{j+2}\) であり、
\[
\delta(ab)=\sigma(a)\delta(b)+\delta(a)b.
\]
この拡大は左 \(B\)-加群として \(\bigoplus_{j\ge0}Bw^j\) であり、
\(\sigma\) が可逆なので右 \(B\)-加群としても
\(\bigoplus_{j\ge0}w^jB\) である。

主結果である定理~\ref{thm:ore} は、
\(\gldim A=n+1\)、\(\ASind A=n+2\) の下で、この Ore 表示と
\(B\) の \(n\) 次元 quadratic Koszul AS-regular 性を与える。
\(n=2,3\) では、必要な AS index を大域次元と最小生成次数から決定できる
（第~\ref{sec:lowdim} 節）。従って、とくに 4 次元 \((1,1,1,2)\) 型は、
3 次元 quadratic AS-regular 代数とその Ore データに帰着する。
その記述の同型関係も明記する。

一般の \(n\) について、最小生成次数と \(\gldim A=n+1\) のみから
\(\ASind A=n+2\) を導く主張は、このノートでは証明しない。
第~\ref{sec:counterexample} 節の反例は大域次元 \(n+2\) を持つため、
大域次元 \(n+1\) に限定した主張への反例ではない。

\section{重み \texorpdfstring{\((1^n,2)\)}{(1,...,1,2)} の構造定理}\label{sec:general}

\begin{theorem}\label{thm:ore}
\(A\) を両側 noetherian AS-regular 代数とし、最小生成次数を
\((1^n,2)\)、\(n\ge1\) とする。さらに
\begin{equation}\label{eq:index-hyp}
d:=\gldim A=n+1,\qquad \ell:=\ASind A=n+2
\end{equation}
を仮定する。このとき \(B=k\langle A_1\rangle\) は両側 noetherian な
\(n\) 次元 quadratic Koszul AS-regular 代数であり、
\[
A\cong B[w;\sigma,\delta],\qquad |w|=2
\]
と書ける。\(\sigma\) は graded automorphism、\(\delta\) は次数 \(2\) の
\(\sigma\)-derivation である。また
\begin{equation}\label{eq:standard-hilbert}
H_B(t)=\frac1{(1-t)^n},\qquad
H_A(t)=\frac1{(1-t)^n(1-t^2)}.
\end{equation}
\end{theorem}

\begin{corollary}\label{cor:hilbert-hyp}
定理~\ref{thm:ore} の AS index の仮定を、
\[
H_A(t)=\frac1{(1-t)^n(1-t^2)}
\]
で置き換えても、同じ結論が成立する。
\end{corollary}

\begin{proof}[系の証明]
AS-Gorenstein 双対性により、\(k_A\) の最小自由分解の最終項は
\(A(-\ell)\) である。第 \(i\) 項（\(0<i<d\)）の shift \(s\) は
\(s\le\ell-(d-i)<\ell\) を満たすので、Euler 和
\(H_A(t)^{-1}\) の最高次数は \(\ell\) である。
仮定の逆 Hilbert 級数の次数は \(n+2\) だから、
\(\ell=n+2\) となる。
ここで用いた最小自由分解の有限性と双対性は
\cite[Lemma~1.1]{LPWZ} による。
\end{proof}

以下、定理~\ref{thm:ore} を証明する。まず最小自由分解の次数を制限し、
\(k[w]\) 上の Tor の有限次元性から交換関係の係数を可逆にする。
次に第 3 微分を用いて、その交換関係が quadratic relations と両立することを示す。

\subsection{最小自由分解と Tor}

\(V=A_1\)、\(\dim_kV=n\) とし、次数 \(2\) の最小生成元を \(w\) とする。
\(k_A\) の最小自由分解を \(P_\bullet\) と書く。
その第 \(i\) 項の shift は、最小性から \(i\) 以上である。
双対分解の最小性と \eqref{eq:index-hyp} から、同じ shift は
\(\ell-(d-i)=i+1\) 以下である。
従って
\begin{equation}\label{eq:two-strands}
P_i=A(-i)^{a_i}\oplus A(-i-1)^{b_i},
\end{equation}
と書け、端の値は
\[
a_0=1,\quad b_0=0,\qquad
a_1=n,\quad b_1=1,\qquad
a_d=0,\quad b_d=1
\]
である。

最小関係式の次数は \(2,3\) のみである。
次数 \(2\) の関係式に \(w\) の線形項があれば \(w\) を消去できるので、
そのような項はない。次数 \(3\) の各単語は \(V\) の元を含む。
従って \(V\) の生成する両側 ideal による商は
\begin{equation}\label{eq:quotient-Q}
Q:=A/(V)\cong k[w].
\end{equation}

\(P_\bullet\otimes_AQ\) において、微分の次数 \(1\) の係数は消え、
次数 \(0\) の係数は最小性から零である。
従って非零になり得る微分は
\begin{equation}\label{eq:tor-block}
Q(-i-1)^{b_i}\xrightarrow{\;wM_i\;}
Q(-i+1)^{a_{i-1}},\qquad
M_i\in M_{a_{i-1},b_i}(k)
\end{equation}
のみである。各直和因子はこのうち 1 本の矢印にだけ関与する。

左 noetherian 性により、左 \(A\)-加群 \(Q\) は各項が有限生成自由加群である
分解を持つ。それを用いて Tor を計算すれば
\(\dim_k\Tor_i^A(k,Q)<\infty\) である。
\(M_i\) に核があれば \eqref{eq:tor-block} の核に自由 \(Q\)-直和因子が生じ、
余核があれば同じ矢印の余核に自由 \(Q\)-直和因子が生じる。
いずれも Tor の有限次元性に反する。よって
\begin{equation}\label{eq:betti-pairing}
b_i=a_{i-1},\qquad M_i\text{ は正方可逆行列}
\end{equation}
である。反対環にもこの議論を適用できる。

\subsection{Hilbert 級数と Betti 数}

\eqref{eq:two-strands} と \eqref{eq:betti-pairing} から
\[
H_A(t)^{-1}=(1-t^2)q(t),\qquad
q(t)=\sum_{i=0}^{n}(-1)^i a_it^i.
\]
ここで \(a_n=b_{n+1}=1\) だから
\[
\deg q=n,\quad q(0)=1,\quad [t]q=-n,\quad [t^n]q=(-1)^n.
\]
有限 GK 次元により \(H_A(t)\) は \(|t|<1\) で収束するため、
\(H_A(t)^{-1}\) はその円板内に根を持たない。
従って \(q\) の根の絶対値はすべて \(1\) 以上である。
定数項と最高次係数の絶対値がともに \(1\) なので根の積の絶対値は \(1\)、
よって全根の絶対値は \(1\) である。
\[
q(t)=\prod_{j=1}^{n}(1-\mu_jt)
\]
と書くと、\(|\mu_j|=1\) かつ \(\sum_j\mu_j=n\) である。
実部を取れば各 \(\operatorname{Re}\mu_j\le1\) の和が \(n\) なので、
すべて \(\mu_j=1\) となる。従って
\begin{equation}\label{eq:betti-binomial}
q(t)=(1-t)^n,\qquad
a_i=\binom ni,\qquad b_i=\binom n{i-1}.
\end{equation}
ここでは範囲外の二項係数を \(0\) とする。
これにより \eqref{eq:standard-hilbert} の \(H_A\) の式を得た。

\subsection{交換関係と quadratic relations の両立}

\eqref{eq:betti-pairing} より最小 cubic relations は \(b_2=n\) 本である。
第 2 微分における \(M_2\) の可逆性は、cubic relations から
\(Vw\) への係数射影が可逆であることを意味する。
反対環での同じ主張により、\(wV\) への係数射影も可逆である。
従って関係式を
\begin{equation}\label{eq:cubic-normal}
wv=\tau(v)w+\partial(v)\qquad(v\in V)
\end{equation}
と正規化できる。ただし \(\tau:V\xrightarrow{\sim}V\)、
\(\partial:V\to T_3(V)\) である。
\(\tau\) を自由代数 \(T=T(V)\) の automorphism に、
\(\partial\) を \(T\) 上の \(\tau\)-derivation に延長する。

\(R\subset V\otimes V\) を quadratic relation space とし、
\(r=\dim_kR=a_2\)、\(I=(R)\subset T\) と置く。
基底を行ベクトル \(\mathbf x=(x_1,\ldots,x_n)\) で表し、
\[
\tau(\mathbf x)=\mathbf xS,\qquad
\partial(\mathbf x)=\mathbf xD,
\qquad S\in GL_n(k),\quad D\in M_n(T_2)
\]
と書く。\(\mathbf xL\) の \(r\) 成分が \(R\) の基底となるように
\(L\in M_{n,r}(V)\) を取る。
右自由加群の元を列ベクトルで表すと、第 2 微分は
\begin{equation}\label{eq:d2-matrix}
d_2=
\begin{pmatrix}L&-Sw-D\\0&\mathbf x\end{pmatrix}.
\end{equation}

\(r=0\) なら \(I=0\) であり、\(\tau,\partial\) は既に \(T/I\) 上に定義される。
以下 \(r>0\) とする。このとき \(P_3\) が存在し、
その次数 \(4\) の直和因子は \(b_3=a_2=r\) 個である。
対応する第 3 微分の列を
\begin{equation}\label{eq:d3-matrix}
d_3\big|_{A(-4)^r}=
\begin{pmatrix}Pw+N\\K\end{pmatrix},
\quad
P\in GL_r(k),\quad N\in M_r(T_2),\quad K\in M_{n,r}(V)
\end{equation}
と書く。\(P\) の可逆性は \eqref{eq:betti-pairing} の \(i=3\) の場合である。

次数 \(3\) の関係式空間は、\(VR+RV\) と選んだ cubic relations の和である。
後者から \(wV\) への射影が可逆なので、\(wV\) を消去すると
\begin{equation}\label{eq:degree-three-splitting}
A_3\cong T_3/(VR+RV)\oplus Vw
\end{equation}
を得る。

\(d_2d_3=0\) の下段は
\begin{equation}\label{eq:relations-xK}
\mathbf xK\in R^{\oplus r}
\end{equation}
を与える。上段で \eqref{eq:cubic-normal} を用い、
\eqref{eq:degree-three-splitting} の成分を比較すると
\begin{align}
LP&=S\tau(K),\label{eq:tau-compat}\\
LN-DK-S\partial(K)&\in M_{n,r}(I_3).
\label{eq:delta-compat}
\end{align}
\(L\) の列は \(k\)-線形独立であり、\(P,S,\tau\) は可逆なので、
\eqref{eq:tau-compat} から \(K\) の列も独立である。
写像 \(V^{\oplus n}\to T_2\)、\((f_i)\mapsto\sum_i x_if_i\) は
ベクトル空間の同型である。従って \eqref{eq:relations-xK} の
\(\mathbf xK\) の成分は \(R\) の基底をなす。

これらの基底に関して
\[
\tau(\mathbf xK)=\mathbf xS\tau(K)=(\mathbf xL)P
\]
なので \(\tau(R)=R\) である。また \eqref{eq:delta-compat} に左から
\(\mathbf x\) を掛けると
\[
\partial(\mathbf xK)
=\mathbf xDK+\mathbf xS\partial(K)
\equiv(\mathbf xL)N\pmod{I_4},
\]
ゆえに \(\partial(R)\subseteq I\) である。
\(\tau\)-derivation の積公式から \(\partial(I)\subseteq I\) も従う。

従って \(\tau,\partial\) は \(\widetilde B=T/I\) 上の
graded automorphism \(\sigma\) と次数 \(2\) の
\(\sigma\)-derivation \(\delta\) に降りる。
\(A\) の全最小関係式は \(R\) と \eqref{eq:cubic-normal} なので、
生成元と関係式の表示を比較して
\begin{equation}\label{eq:ore-presentation}
A\cong \widetilde B[w;\sigma,\delta]
\end{equation}
を得る。Ore 拡大の自由加群表示により \(\widetilde B\) は \(A\) に埋め込まれ、
その像は \(k\langle A_1\rangle=B\) である。
また \eqref{eq:standard-hilbert} の \(H_A\) の式と
\(H_A=H_B/(1-t^2)\) から、\(H_B\) の式も従う。

\subsection{AS-regular 性の降下}\label{subsec:descent}

残る \(B\) の AS-regular 性を示す。
\(A\) が左右の \(B\)-自由加群であることから、右 ideal \(J\subset B\) に対して
\(JA\cap B=J\)、左 ideal に対して \(AJ\cap B=J\) が成り立つ。
従って ideal の昇鎖を延長・縮約することにより、\(B\) も両側 noetherian である。

\(B_+=B_{\ge1}\) と置く。\(\sigma(B_+)=B_+\)、
\(\delta(B_+)\subseteq B_+\) なので、\(B_+A\) は両側 ideal である。
右 \(A\)-加群として
\[
Q=k\otimes_BA=A/B_+A\cong k[w],
\qquad Q_B\cong\bigoplus_{j\ge0}k(-2j).
\]
\(Q\) の長さ高々 \(d\) の \(A\)-projective resolution を \(B\) に制限する。
\(A_B\) は自由なので、これにより \(\pd_BQ\le d\) を得る。
\(k_B\) は \(Q_B\) の直和因子であるから \(\pd_Bk\le d\) である。

\(B\) の noetherian 性により、\(k_B\) は有限生成自由加群からなる有限分解を持つ。
この分解を用いた誘導・制限の随伴と \(A_B\cong\bigoplus_{j\ge0}B(-2j)\) から、
次数付きベクトル空間として
\begin{equation}\label{eq:ext-induction}
E^i:=\underline{\Ext}^{i}_A(Q,A)
\cong\underline{\Ext}^{i}_B(k,A_B)
\cong\bigoplus_{j\ge0}
\underline{\Ext}^{i}_B(k,B)(-2j).
\end{equation}
各 \(E^i\) は各次数有限次元で次数下有界である。

右 \(A\)-加群の短完全列
\[
0\longrightarrow Q(-2)\xrightarrow{\cdot w}Q
\longrightarrow k\longrightarrow0
\]
に \(\underline{\Ext}_A(-,A)\) を適用する。
AS-Gorenstein 条件により、\(i<d-1\) では \(E^i\cong E^i(2)\) であり、
\(E^d\to E^d(2)\) は全射である。
最低の非零次数を考えれば、下有界性によりこれらの \(E^i\) はすべて零となる。
また \(i>d\) では大域次元により零である。残る部分は
\[
0\longrightarrow E^{d-1}\longrightarrow E^{d-1}(2)
\longrightarrow k(\ell)\longrightarrow0
\]
なので
\[
H_{E^{d-1}}(t)=\frac{t^{2-\ell}}{1-t^2}.
\]
これと \eqref{eq:ext-induction} を比較すると
\begin{equation}\label{eq:ext-B}
\underline{\Ext}^{i}_B(k,B)\cong
\begin{cases}k(\ell-2)=k(n)&i=d-1=n,\\0&i\ne n\end{cases}
\end{equation}
を得る。非零の Ext は 1 次元で 1 つの内部次数に集中するため、
\(B_+\) の作用は零であり、上式は左 \(B\)-加群の同型でもある。

有限最小自由分解の最高次の Ext は非零である。
実際、最終の双対微分の像は自由加群の \(B_+\) による部分加群に入るため、
その余核は零にならない。従って \eqref{eq:ext-B} から \(\pd_Bk=n\) である。
連結次数付き代数に対する \(\gldim B=\pd_Bk\)
（\cite[Proposition~2.1 の証明]{CKS} も参照）から大域次元は \(n\) である。

反対環は \(B^{\mathrm{op}}\) 上の Ore 拡大で、automorphism は
\(\sigma^{-1}\)、derivation は \(-\delta\sigma^{-1}\) である。
従って同じ議論を反対環に適用し、左側の AS-Gorenstein 条件も得る。
有限 GK 次元は \(B\subseteq A\) から従う。
これで \(B\) は \(n\) 次元 AS-regular、AS index \(n\) と分かった。
その最小分解の第 \(i\) 項の shift は、最小性と双対性により \(i\) に限られる。
従って \(B\) は Koszul である。\(B=T(V)/(R)\) なので quadratic でもあり、
定理~\ref{thm:ore} の証明が完了した。

\section{低次元の場合と次数付き同型}
\label{sec:lowdim}

大域次元が3または4の場合には、最小生成元の次数だけから
定理\ref{thm:ore}に必要なAS indexを決定できる。
以下では両側 noetherian性から得られる成長度の下界と、最小自由分解の双対性を用いる。

\subsection{最小自由分解とHilbert級数の零点}

$A$を大域次元$d$、AS index$\ell$の連結次数付きAS-regular代数とし、
$k_A=A/A_{\geq 1}$の最小次数付き右自由分解を
\[
  0\longrightarrow P_d\longrightarrow\cdots\longrightarrow P_0
  \longrightarrow k_A\longrightarrow0,
  \qquad
  P_i=\bigoplus_{\nu=1}^{b_i}A(-a_{i\nu})
\]
と書く。ここで$A(-a)_m=A_{m-a}$である。
AS双対性によれば、$P_0=A$、$P_d=A(-\ell)$であり、
$P_{d-i}$のシフトの多重集合は
$\{\ell-a_{i\nu}\}_{\nu=1}^{b_i}$に一致する
\cite[Lemma 1.1]{LPWZ}。
また
\begin{equation}
  a_{i\nu}\geq i
  \label{eq:lowdim-shift-bound}
\end{equation}
である。実際、最小性より$P_i\otimes_A k\cong\Tor_i^A(k,k)$であり、
$\Tor_i^A(k,k)$を計算する被約bar複体の第$i$項
$A_{\geq1}^{\otimes i}$は内部次数$i$未満の成分を持たない。

自由分解の各次数でEuler標数を取ると
\begin{equation}
  H_A(t)=\frac{1}{p_A(t)},
  \qquad
  p_A(t)=\sum_{i=0}^{d}(-1)^i\sum_{\nu=1}^{b_i}t^{a_{i\nu}}
  \label{eq:lowdim-euler}
\end{equation}
を得る。以下で用いる代数は次数1と2の元で生成される。
この場合、$t=1$における$p_A(t)$の零点の位数を$m$とすると
$\GKdim A\leq m$である。ここでは$p_A(1)\ne0$のとき$m=0$とする。
この不等式は、$h_j=\dim_k A_j$と置いたときの
\[
  \sum_{j=0}^{N}h_j
  \leq eH_A(e^{-1/N})=O(N^m)
\]
から従う。AS-regular性に含まれる有限GK次元の条件により、
この式のHilbert級数は$0<t<1$で収束する。
また、$1$と固定した最小生成元が張る空間を$U$とすれば
$U^N\subseteq\bigoplus_{j=0}^{2N}A_j$なので、
上の評価はGK次元の定義に適用できる。
両側 noetherianかつ$d\geq3$ならば
\cite[Lemma 1.2]{LPWZ}より$\GKdim A\geq3$であり、したがって
\begin{equation}
  p_A(1)=p_A'(1)=p_A''(1)=0
  \label{eq:lowdim-three-zeros}
\end{equation}
となる。

\subsection{生成元の次数が\texorpdfstring{$(1,1,1,2)$}{(1,1,1,2)}の場合}

\begin{proposition}
\label{prop:index-four}
$A$を両側 noetherianな連結次数付きAS-regular$\C$代数とする。
$\gldim A=4$であり、最小生成元の次数が$(1,1,1,2)$ならば、
\[
  \ASind A=5,
  \qquad
  H_A(t)=\frac{1}{(1-t)^3(1-t^2)}.
\]
さらに、$B=\C\langle A_1\rangle\subseteq A$は
3次元のquadratic Koszul AS-regular代数であり、ある
$\sigma\in\Aut_{\mathrm{gr}}(B)$と次数2の$\sigma$導分$\delta$に対して
\[
  A\cong B[w;\sigma,\delta],\qquad \deg w=2
\]
となる。この同型は$B$上の恒等写像を延長する。
\end{proposition}

\begin{proof}
AS双対性で最小分解の両端を決め、
\eqref{eq:lowdim-three-zeros}から中間項のシフトの平方和を求める。
各シフトが入る区間上の不等式がAS indexの上界を与える。

AS indexを$\ell$と書く。最小生成元の次数とAS双対性より
\[
\begin{split}
  P_1&=A(-1)^{\oplus3}\oplus A(-2),\\
  P_3&=A(-\ell+1)^{\oplus3}\oplus A(-\ell+2),
  \qquad P_4=A(-\ell).
\end{split}
\]
$P_3$にはシフト$\ell-2$の直和因子があるので、
\eqref{eq:lowdim-shift-bound}から$\ell\geq5$を得る。
$P_2=\bigoplus_{j=1}^{r}A(-s_j)$と書けば、同じ評価と双対性により
\begin{equation}
  2\leq s_j\leq\ell-2,
  \qquad
  \{s_1,\ldots,s_r\}=\{\ell-s_1,\ldots,\ell-s_r\}
  \label{eq:four-middle-shifts}
\end{equation}
が多重集合として成立する。
\eqref{eq:lowdim-euler}の分母は
\[
  p_A(t)=1-3t-t^2+\sum_{j=1}^{r}t^{s_j}
       -3t^{\ell-1}-t^{\ell-2}+t^\ell.
\]
$p_A(1)=0$から$r=6$であり、
\eqref{eq:four-middle-shifts}の対称性から
\begin{equation}
  \sum_{j=1}^{6}s_j=3\ell
  \label{eq:four-first-moment}
\end{equation}
を得る。さらに$p_A''(1)=0$は
\[
  -2+\sum_{j=1}^{6}s_j(s_j-1)
  -3(\ell-1)(\ell-2)-(\ell-2)(\ell-3)+\ell(\ell-1)=0
\]
であるから、\eqref{eq:four-first-moment}を代入して
\begin{equation}
  \sum_{j=1}^{6}s_j^2=3\ell^2-10\ell+14
  \label{eq:four-second-moment}
\end{equation}
を得る。一方、$2\leq s_j\leq\ell-2$より
\[
  s_j(\ell-s_j)-2(\ell-2)
   =(s_j-2)(\ell-2-s_j)\geq0.
\]
これを$j$について加え、
\eqref{eq:four-first-moment}と\eqref{eq:four-second-moment}を使うと
\[
  12\ell-24\leq
  \sum_{j=1}^{6}s_j(\ell-s_j)=10\ell-14.
\]
したがって$\ell\leq5$であり、先の下界と合わせて$\ell=5$となる。
このとき$s_j\in\{2,3\}$かつ$\sum_js_j=15$なので、
次数2と次数3の最小関係式はそれぞれ3本である。
よって
\[
  p_A(t)=1-3t+2t^2+2t^3-3t^4+t^5
        =(1-t)^3(1-t^2).
\]

定理\ref{thm:ore}を$n=3$に適用すると、
$B=\C\langle A_1\rangle$は3次元のquadratic Koszul AS-regular代数であり、
主張のOre表示が得られる。
\end{proof}

\subsection{生成元の次数が\texorpdfstring{$(1,1,2)$}{(1,1,2)}の場合}

\begin{proposition}
\label{prop:index-three}
$A$を両側 noetherianな連結次数付きAS-regular$\C$代数とする。
$\gldim A=3$であり、最小生成元の次数が$(1,1,2)$ならば、
\[
  \ASind A=4,
  \qquad
  H_A(t)=\frac{1}{(1-t)^2(1-t^2)}.
\]
さらに$B=\C\langle A_1\rangle$は、次数1の生成元$x,y$を適切に選ぶと
\begin{equation}
  B_q=\C\langle x,y\rangle/(yx-qxy)\quad(q\in\C^\times),
  \qquad\text{または}\qquad
  B_J=\C\langle x,y\rangle/(yx-xy-x^2)
  \label{eq:plane-normal-forms}
\end{equation}
のいずれかである\cite[\S 2, Corollary 2.2の後]{RZ}。
ある$\sigma\in\Aut_{\mathrm{gr}}(B)$と次数2の$\sigma$導分$\delta$に対して、
$B$上の恒等写像を延長する次数付き同型
\[
  A\cong B[w;\sigma,\delta],\qquad\deg w=2
\]
が存在する。
\end{proposition}

\begin{proof}
AS indexを$\ell$と書く。最小分解の双対性から
\[
  p_A(t)=1-2t-t^2+2t^{\ell-1}+t^{\ell-2}-t^\ell
\]
である。\eqref{eq:lowdim-three-zeros}のうち$p_A'(1)=0$を用いると
\[
  0=-2-2+2(\ell-1)+(\ell-2)-\ell=2\ell-8
\]
となるので、$\ell=4$である。したがって
\[
  p_A(t)=1-2t+2t^3-t^4=(1-t)^2(1-t^2).
\]
定理\ref{thm:ore}を$n=2$に適用すると、
$B=\C\langle A_1\rangle$は2次元AS-regular代数であり、
主張のOre表示と$H_B(t)=(1-t)^{-2}$が得られる。

ここで$B$の表示を確認する。
$B$は次数1の2元$x,y$で最小生成されるので、
AS双対性と$H_B(t)=(1-t)^{-2}$から最小右自由分解は
\[
  0\longrightarrow B(-2)\xrightarrow{d_2}B(-1)^{\oplus2}
   \xrightarrow{d_1}B\longrightarrow k_B\longrightarrow0
\]
となる。$d_1(e_1)=x$、$d_1(e_2)=y$とし、
$d_2(1)=e_1a+e_2b$と書くと、$a,b\in B_1$であり、
$B$の唯一の最小関係式は
$r=x\otimes a+y\otimes b\in B_1\otimes B_1$である。
この分解に$\operatorname{Hom}_B(-,B)$を施すと、AS条件から
\[
  B/(Ba+Bb)\cong k
\]
が得られる。従って$a,b$は$B_1$の基底であり、
$r$の係数行列は可逆である。

$r$の可換化を$\overline r\in\operatorname{Sym}^2(B_1)$とする。
$\overline r=0$なら$r$は$xy-yx$の非零スカラー倍であり、
$B\cong\C[x,y]$、すなわち$q=1$の量子平面となる。
$\overline r$の階数が2ならば、$\C$上の変数変換で
$\overline r$を$xy$の非零スカラー倍にできる。
この座標では$r=\alpha xy+\beta yx$となり、
係数行列の可逆性により$\alpha\beta\ne0$だから、
$q=-\alpha/\beta$と置けば量子平面の表示を得る。
$\overline r$の階数が1なら、スカラー倍と変数変換により
\[
  r=x^2+\gamma(xy-yx)
\]
と書ける。係数行列が可逆なので$\gamma\ne0$であり、
$y'=\gamma y$と置くと$y'x-xy'-x^2=0$となり、Jordan平面の表示を得る。
これで\eqref{eq:plane-normal-forms}の全場合を得た。
\end{proof}

\begin{corollary}
\label{cor:lowdim-converse}
$n=2$の場合は$B$を\eqref{eq:plane-normal-forms}のいずれかとし、
$n=3$の場合は$B$を次数1の3元で生成される両側 noetherianな
quadratic Koszul AS-regular代数で$\gldim B=3$を満たすものとする。
$\sigma\in\Aut_{\mathrm{gr}}(B)$と次数2の$\sigma$導分$\delta$を任意に取る。
このとき$B[w;\sigma,\delta]$は両側 noetherianなAS-regular代数であり、
その大域次元は$n+1$、最小生成元の次数は$(1^n,2)$である。
\end{corollary}

\begin{proof}
量子平面とJordan平面は、それぞれ
$\C[x][y;\tau]$、$\C[x][y;\partial]$というOre拡大である。
ここで$\tau(x)=qx$、$\partial(x)=x^2$であり、$\deg x=\deg y=1$とする。
これらは両側 noetherianな2次元AS-regular代数であり、Hilbert級数は$(1-t)^{-2}$である。
これらの性質と、一般の$A=B[w;\sigma,\delta]$のAS-regular性には、
次のOre拡大定理を適用できる。
noetherian AS-regularな連結次数付き代数はtwisted Calabi--Yauであり
\cite[\S 1.2]{LWW}、自己同型によるOre拡大はその次元を1増やす
\cite[Theorem 3.3]{LWW}。
得られた代数のAS条件と有限大域次元は\cite[Lemma 1.3]{LWW}から従い、
両側 noetherian性はOre拡大に対するHilbertの基底定理による
\cite[\S 1.3]{LWW}。
これらを$\C[x]$からの拡大にも適用すれば、最初に述べた2つの平面の性質も得られる。

いずれの$n$についても$H_B(t)=(1-t)^{-n}$である。
$n=3$のときは、Koszul性とAS双対性による線形分解
$0\to B(-3)\to B(-2)^3\to B(-1)^3\to B\to k\to0$から従う。
$A$の左$B$自由基底が$\{w^j\}_{j\geq0}$なので
\[
  H_A(t)=\frac{H_B(t)}{1-t^2}
        =\frac{1}{(1-t)^n(1-t^2)},
\]
従って有限GK次元の条件も満たされる。
$A_1=B_1$、$A_2=B_2\oplus\C w$であるから、
$B$の最小生成元に$w$を加えたものは最小生成系である。
\end{proof}

\subsection{Oreデータの同値関係}

次数1で生成される部分代数$\C\langle A_1\rangle$は
次数付き同型によって保たれる。
したがって、上述のOre表示の同型判定は次の形に尽きる。

\begin{proposition}
\label{prop:ore-isomorphism}
$B,B'$を次数1で生成される連結次数付き$\C$代数とし、
\[
  A=B[w;\sigma,\delta],\qquad
  A'=B'[w';\sigma',\delta'],\qquad
  \deg w=\deg w'=2
\]
を次数付きOre拡大とする。
ここで$\sigma,\sigma'$は次数付き自己同型、
$\delta,\delta'$は次数2の対応する歪導分である。
次数付き$\C$代数同型$F:A\to A'$は、
次のデータと1対1に対応する：
次数付き同型$\phi:B\to B'$、$\lambda\in\C^\times$、$u\in B'_2$であって
\begin{align}
  \phi\circ\sigma&=\sigma'\circ\phi,
  \label{eq:ore-iso-sigma}\\
  \phi\bigl(\delta(b)\bigr)
    &=\lambda\delta'\bigl(\phi(b)\bigr)
      +u\phi(b)-\sigma'\bigl(\phi(b)\bigr)u
      \qquad(b\in B)
  \label{eq:ore-iso-delta}
\end{align}
を満たすものである。
この対応は
\[
  F|_B=\phi,\qquad F(w)=\lambda w'+u
\]
で与えられる。
\end{proposition}

\begin{proof}
$A_1=B_1$、$A'_1=B'_1$なので、次数付き同型$F$は
$\C\langle A_1\rangle=B$を$B'$へ写し、次数付き同型$\phi:B\to B'$を誘導する。
$A_2=B_2\oplus\C w$、$A'_2=B'_2\oplus\C w'$であるから、
ある一意な$\lambda\in\C^\times$と$u\in B'_2$により
$F(w)=\lambda w'+u$と書ける。
$wb=\sigma(b)w+\delta(b)$に$F$を施し、左$B'$自由基底の$w'$係数と定数項を比較すると
\begin{align*}
  (\lambda w'+u)\phi(b)
   &=\lambda\sigma'\bigl(\phi(b)\bigr)w'
     +\lambda\delta'\bigl(\phi(b)\bigr)+u\phi(b),\\
  \phi\bigl(\sigma(b)\bigr)(\lambda w'+u)+\phi\bigl(\delta(b)\bigr)
   &=\lambda\phi\bigl(\sigma(b)\bigr)w'
     +\phi\bigl(\sigma(b)\bigr)u+\phi\bigl(\delta(b)\bigr)
\end{align*}
より\eqref{eq:ore-iso-sigma}と\eqref{eq:ore-iso-delta}を得る。

逆にこれらの等式が成り立てば、$F|_B=\phi$、$F(w)=\lambda w'+u$は
Ore関係式を保つので、次数付き準同型$F:A\to A'$を定める。
その像は$B'$および$w'=\lambda^{-1}(F(w)-u)$を含むため、$F$は全射である。
また、$b_m\ne0$である元$\sum_{i=0}^{m}b_iw^i$の像を
左$B'$自由基底で展開したとき、$(w')^m$の係数は
$\lambda^m\phi(b_m)\ne0$である。
従って$F$は単射でもある。
\end{proof}

\eqref{eq:ore-iso-sigma}と\eqref{eq:ore-iso-delta}は$B_1$上で確認すれば十分である。
前者の両辺は代数準同型であり、前者が成立すると、後者の両辺は
$B$から$B'$への同じ$\phi\circ\sigma$および$\phi$に関する歪導分になるためである。
具体的には、その差を$D$とすると
$D(ab)=\phi(\sigma(a))D(b)+D(a)\phi(b)$であり、
$D|_{B_1}=0$から積の長さに関する帰納法で$D=0$を得る。
特に同じ$B$の上で変数を$\widetilde w=\lambda w+u$（$u\in B_2$）に取り替えると、
$\sigma$は変わらず、歪導分は
\[
  \widetilde\delta(b)=\lambda\delta(b)+ub-\sigma(b)u
  \qquad(b\in B)
\]
に変わる。

\section{\texorpdfstring{$(1,1,2)$}{(1,1,2)} 型代数の次数 \texorpdfstring{$1$}{1} 中心拡大}
\label{sec:central}

代数 $C$ の\emph{次数 $1$ の正則中心拡大}とは、
$t\in Z(A)_1$ が非零因子であり、次数付き同型 $A/(t)\cong C$ を持つ代数 $A$ をいう。
ここで $Z(A)$ は $A$ の中心である。Hilbert 級数の変数には $s$ を用いる。

\subsection{中心拡大から Ore データへ}

まず、後で反復にも用いる中心拡大の性質をまとめる。

\begin{lemma}\label{lem:central-lift}
$D$ を連結局所有限次数付き代数、$u\in Z(D)_e$ を非零因子とし、
$e\ge1$、$E=D/(u)$ とおく。
$E$ が両側 noetherian な $d$ 次元 AS-regular代数で AS index $\ell$ を持つなら、
$D$ は両側 noetherian な $d+1$ 次元 AS-regular代数であり、
\[
  H_D(s)=\frac{H_E(s)}{1-s^e},\qquad
  \ASind(D)=\ell+e
\]
が成り立つ。さらに、$e=1$ で $E$ の最小生成次数が $(1^m,2)$ なら、
$D$ の最小生成次数は $(1^{m+1},2)$ である。
\end{lemma}

\begin{proof}
正則中心拡大の議論 \cite[Propositions~2.1 and~2.3]{CKS} を、
本稿の重み付き次数に合わせて示す。完全列
\[
  0\longrightarrow D(-e)\xrightarrow{\cdot u}D
    \longrightarrow E\longrightarrow0
\]
から Hilbert 級数の等式を得る。従って $E$ の多項式増大度から $D$ の
有限 GK 次元も従う。右 $D$-加群として $\pd_D E=1$ である。
長さ $d$ の $E$-自由分解で $k$ を分解すると、各項の $D$ 上の射影次元は
高々 $1$ なので、短完全列に分けた次元移動により $\pd_D k\le d+1$ となる。
上の完全列に $\underline{\Ext}_D^*(-,D)$ を適用すると、
\[
  \underline{\Ext}_D^q(E,D)=
  \begin{cases}
    E(e)&q=1,\\
    0&q\ne1.
  \end{cases}
\]
したがって環変更のスペクトル系列は $q=1$ の行に退化し、
\begin{equation}\label{eq:central-ext-shift}
  \underline{\Ext}_D^{i+1}(k,D)
  \cong\underline{\Ext}_E^i(k,E)(e)
\end{equation}
を与える。$i=d$ の右辺は $k(\ell+e)\ne0$ だから $\pd_D k=d+1$ であり、
連結次数付き代数に対する $\gldim D=\pd_D k$ により大域次元も $d+1$ である。
左加群についても同じ計算が成立し、左右の AS 条件と AS index $\ell+e$ を得る。

Noetherian 性も正の次数だけを用いて持ち上げられる。
減少 filtration $F^pD=u^pD$ に対して、中心性と正則性から
$\operatorname{gr}_uD\cong E[U]$ である。ここで $U$ は中心変数で、
filtration 次数 $1$、内部次数 $e$ を持つ。
同次右 ideal $J\subset D$ に誘導 filtration を入れると、
$\operatorname{gr}_uJ$ は noetherian 環 $E[U]$ の同次右 ideal である。
その有限個の二重同次生成元を $J$ へ持ち上げる。
内部次数 $N$ の元をこれらの元で近似し、その差の $u$-adic order を
順次増やすと、order は $N/e$ 以下なので補正は有限回で停止する。
従って $J$ はそれらの持ち上げで生成される。
任意の右 ideal に対しては内部次数による増加 filtration を用い、
同次 ideal の生成元を持ち上げて最高次数を減らせば有限生成性が従う。
左 ideal にも同じ議論が適用できるので、$D$ は両側 noetherian である。

$e=1$ の場合を考える。
$E$ の最小生成元を同次的に $D$ へ持ち上げ、それらに $u$ を加える。
これらが $D$ を生成することは次数に関する帰納法で従う。実際、
同次元の $E$ での像を生成元の式として持ち上げると、差は $uD$ に属し、
$u$ を除いた係数の次数は $1$ 小さい。また $D_1\to E_1$ の核は $ku$ である。
次数 $2$ の生成元の持ち上げが $D_1$ の積で書けるなら、その像も $E_1$ の積で書け、
$E$ における最小性に反する。よって指定された最小生成次数を得る。
\end{proof}

\begin{theorem}\label{thm:central-ore}
$C$ を両側 noetherian な $3$ 次元 AS-regular代数で、
最小生成次数が $(1,1,2)$ であるものとする。
$A/(t)\cong C$ を次数 $1$ の正則中心拡大とする。
このとき $B=\langle A_1\rangle$ は両側 noetherian なquadratic Koszul AS-regular代数で、
$\gldim B=3$ である。$A_2/(A_1A_1)$ の基底を持ち上げる元 $w\in A_2$ を選ぶと、
ある $\sigma\in\Aut_{\mathrm{gr}}(B)$ と次数 $2$ の $\sigma$ 微分 $\delta$ により
\begin{equation}\label{eq:central-ore}
 A\cong B[w;\sigma,\delta],\qquad
 t\in Z(B)_1,\qquad \sigma(t)=t,\qquad\delta(t)=0
\end{equation}
と書ける。さらに、$R=B/(t)$ は量子平面または Jordan 平面であり、
商への誘導写像を上線で表すと
\begin{equation}\label{eq:central-quotient-ore}
 C\cong R[w;\bar\sigma,\bar\delta].
\end{equation}
その Hilbert 級数は
\[
 H_B(s)=\frac1{(1-s)^3},\qquad
 H_C(s)=\frac1{(1-s)^2(1-s^2)},\qquad
 H_A(s)=\frac1{(1-s)^3(1-s^2)}
\]
であり、$\gldim A=4$、$\ASind(A)=5$ である。

逆に、$B$ を両側 noetherian な二次 $3$ 次元 AS-regular代数とし、
$t\in Z(B)_1$ を非零因子で、$R=B/(t)$ が $2$ 次元 AS-regularとなるものとする。
$\sigma\in\Aut_{\mathrm{gr}}(B)$ と次数 $2$ の $\sigma$ 微分 $\delta$ が
$\sigma(t)=t$、$\delta(t)=0$ を満たすなら、$B[w;\sigma,\delta]$ は
\eqref{eq:central-quotient-ore} で定まる $(1,1,2)$ 型代数の次数 $1$ の正則中心拡大である。
\end{theorem}

\begin{proof}
主定理を $A$ と $C$ の双方に適用し、Ore 基底によって
$B/(t)$ を $\langle C_1\rangle$ と同定する。これにより、中心元の条件と商の AS-regular性を得る。

命題~\ref{prop:index-three} により $\ASind(C)=4$ である。
補題~\ref{lem:central-lift} から、$A$ は両側 noetherian な $4$ 次元 AS-regular代数で、
最小生成次数は $(1,1,1,2)$、AS indexは $5$ である。
定理~\ref{thm:ore} を $n=3$ で適用すると、
$B=\langle A_1\rangle$ はquadratic Koszul $3$ 次元 AS-regular代数となり、
$A=B[w;\sigma,\delta]$ を得る。

$t\in A_1=B_1$ かつ $t\in Z(A)$ なので $t\in Z(B)$ である。
$wt=tw$ と Ore 関係を比較すると
\[
   (\sigma(t)-t)w+\delta(t)=0.
\]
$A=\bigoplus_{j\ge0}Bw^j$ における係数の一意性から、
$\sigma(t)=t$ および $\delta(t)=0$ が従う。
したがって $\sigma$ と $\delta$ は $(t)\subset B$ を保ち、$R=B/(t)$ に
$\bar\sigma\in\Aut_{\mathrm{gr}}(R)$ と次数 $2$ の $\bar\sigma$ 微分 $\bar\delta$ を誘導する。
中心性と Ore 基底から
\[
  tA=\bigoplus_{j\ge0}tBw^j,\qquad
  B\cap tA=tB,\qquad
  A/tA=\bigoplus_{j\ge0}(B/tB)w^j
\]
となる。商における積は誘導された Ore 関係で与えられるので、
\eqref{eq:central-quotient-ore} は実際に次数付き代数同型である。

一方、$C$ に定理~\ref{thm:ore} を $n=2$ で適用すると、
$\langle C_1\rangle$ は二次 $2$ 次元 AS-regular代数である。
$B$ の $C$ における像はちょうど $\langle C_1\rangle$ であり、その核は上で示した $tB$ だから、
\[
 R=B/(t)\cong\langle C_1\rangle.
\]
$k=\C$ 上の二次 $2$ 次元 AS-regular代数の関係式は、基底変換によって
\begin{equation}\label{eq:central-plane}
  R\cong k\langle x,y\rangle/(yx-qxy-\varepsilon x^2),
  \qquad
  (q,\varepsilon)\in
  \{(q,0)\mid q\in k^\times\}\cup\{(1,1)\}
\end{equation}
の形になる（\cite[Section~2, following Corollary~2.2]{RZ}）。
これはそれぞれ量子平面と Jordan 平面であり、$H_R(s)=(1-s)^{-2}$ である。
$t$ は $B$ 上でも非零因子であるから $H_B=H_R/(1-s)$、
さらに Ore 基底から $H_A=H_B/(1-s^2)$ および $H_C=H_R/(1-s^2)$ を得る。

逆向きでは、Ore 拡大 $A=B[w;\sigma,\delta]$ の左 $B$ 基底を用いる。
$\sigma(t)=t$ と $\delta(t)=0$ より $wt=tw$ なので、$t$ は $A$ の中心に属する。
$t\sum_jb_jw^j=0$ なら各 $tb_j=0$ であるから各 $b_j=0$ となり、
$t$ は $A$ 上でも非零因子である。商を計算する上の議論もそのまま適用できる。
自己同型を用いる次数付き Ore 拡大は、両側 noetherian 性と AS-regular性を保ち、
大域次元を $1$ 増やす
（\cite[Sections~1.2--1.3 and Theorem~3.3]{LWW}）。
これを $B$ と $R$ に適用すると、$A$ と $C=R[w;\bar\sigma,\bar\delta]$ は
それぞれ $4$ 次元と $3$ 次元の AS-regular代数となる。
有限 GK 次元は Ore 基底による上記の Hilbert 級数からも確認できる。
また $C_2=R_2\oplus kw$ なので $w\notin\langle C_1\rangle$ であり、
$C$ の最小生成次数は $(1,1,2)$ である。
\end{proof}

\subsection{次数 \texorpdfstring{$1$}{1} 部分の明示的表示}

\begin{proposition}\label{prop:central-base-presentation}
定理~\ref{thm:central-ore} に現れる $(B,t)$ は、$x,y\in B_1$ の選択により
\begin{equation}\label{eq:central-base-presentation}
 B\cong
 \frac{k\langle t,x,y\rangle}
 {\bigl(tx-xt,\ ty-yt,\ yx-qxy-\varepsilon x^2-t(ax+by)-ct^2\bigr)}
\end{equation}
と表示できる。ここで $\deg t=\deg x=\deg y=1$、$a,b,c\in k$ であり、
$(q,\varepsilon)$ は \eqref{eq:central-plane} の条件を満たす。
逆に、これらの任意のパラメータに対して、\eqref{eq:central-base-presentation} は
両側 noetherian な二次 $3$ 次元 AS-regular整域を定める。
$t$ は正則中心元であり、その商は \eqref{eq:central-plane} の $R$ である。
\end{proposition}

\begin{proof}
まず表示の右辺を $P$ とおく。可換多項式環 $S=k[t,x]$ 上で
\begin{align*}
  &\alpha(t)=t,\qquad &&\alpha(x)=qx+bt,\\
  &\beta(t)=0,\qquad &&\beta(x)=\varepsilon x^2+atx+ct^2
\end{align*}
と定める。$q\ne0$ なので $\alpha$ は次数付き自己同型であり、
$\alpha^{-1}(x)=q^{-1}(x-bt)$ である。
$\beta$ を自由代数から $\alpha$ 微分として延長すると、
$f=\varepsilon x^2+atx+ct^2$ に対して
\[
  \beta(xt-tx)=ft-tf=0\quad\text{in }S
\]
であるから、$S$ 上の $\alpha$ 微分を定める。その次数は $1$ である。
Ore 関係 $ys=\alpha(s)y+\beta(s)$ を $s=t,x$ について書くと
\eqref{eq:central-base-presentation} の関係式に一致するため、
\begin{equation}\label{eq:central-base-iterated-ore}
 P\cong k[t,x][y;\alpha,\beta].
\end{equation}
従って
\[
   \{t^i x^j y^h\mid i,j,h\ge0\}
\]
は $P$ の $k$ 基底であり、$H_P(s)=(1-s)^{-3}$ である。
さらに、$y$ の最高次の係数がそれぞれ $r_m,s_n\in S\setminus\{0\}$ である
非零元の積の最高次係数は $r_m\alpha^m(s_n)$ なので、$P$ は整域である。
$S$ は両側 noetherian な $2$ 次元 AS-regular代数であり、$\alpha$ は自己同型だから、
\cite[Sections~1.2--1.3 and Theorem~3.3]{LWW} により
$P$ は両側 noetherian な $3$ 次元 AS-regular代数である。
有限 GK 次元は既に計算した Hilbert 級数で確認される。
また、関係式から $t$ は中心元であり、PBW 基底からその乗法は単射である。
$t=0$ とすれば商は \eqref{eq:central-plane} に一致する。

次に、定理~\ref{thm:central-ore} の $B$ に戻る。
$R_1$ の基底を \eqref{eq:central-plane} の形に選び、$x,y\in B_1$ へ持ち上げる。
このとき $t,x,y$ は $B_1$ の基底である。
$yx-qxy-\varepsilon x^2$ は $B_2$ に属し、$R$ における像が $0$ なので、
ある $a,b,c\in k$ によって
\[
  yx-qxy-\varepsilon x^2=t(ax+by+ct)
\]
と書ける。したがって同じパラメータの $P$ から $B$ への次数付き全射がある。
両者の Hilbert 級数は $(1-s)^{-3}$ であるため、各次数で全射は同型となり、
$P\cong B$ が従う。
\end{proof}

\subsection{同型によるデータの同一視}

\eqref{eq:central-base-presentation} のパラメータと、その $B$ 上の
$(\sigma,\delta)$ は、生成元の選択に依存する。この依存は次の同値関係で記述できる。
中心直線を指定した拡大のデータを $(B,kt,\sigma,\delta)$ とし、
これに $A=B[w;\sigma,\delta]$ を対応させる。
別のデータを $(B',kt',\sigma',\delta')$、対応する代数を
$A'=B'[w';\sigma',\delta']$ とする。

\begin{proposition}\label{prop:central-isomorphisms}
中心直線を保つ次数付き同型 $F:A\to A'$、$F(kt)=kt'$ は、
次のデータと $1$ 対 $1$ に対応する：
次数付き同型 $\phi:B\to B'$、$\lambda,\mu\in k^\times$ および $h\in B'_2$ であって、
\begin{align}
  \phi(t)&=\mu t',\label{eq:central-iso-t}\\
  \phi\sigma&=\sigma'\phi,\label{eq:central-iso-sigma}\\
  \phi\delta(b)&=\lambda\delta'\phi(b)+h\phi(b)-\phi\sigma(b)h
       \qquad(b\in B)\label{eq:central-iso-delta}
\end{align}
を満たすものである。この対応における同型は
\[
  F|_B=\phi,\qquad F(w)=\lambda w'+h
\]
で与えられる。
\end{proposition}

\begin{proof}
命題~\ref{prop:ore-isomorphism} により、同型 $F$ の形と
\eqref{eq:central-iso-sigma}、\eqref{eq:central-iso-delta} の必要十分性を得る。
$t\in B_1$ なので、$F(kt)=kt'$ という条件はちょうど
\eqref{eq:central-iso-t} である。
\end{proof}

従って次数 $1$ の中心直線を指定した拡大の同型類は、
定理~\ref{thm:central-ore} のデータを
\eqref{eq:central-iso-t}--\eqref{eq:central-iso-delta} で同一視したものと一致する。
中心元そのものを指定する場合は $\mu=1$ とする。
固定した代数 $C$ への商写像まで指定する場合は、
$A/(t)\cong C$ という同型をデータに加え、誘導される商の同型がこの同一視と可換になることを課す。
中心直線を指定しない代数の同型を考える場合は、\eqref{eq:central-iso-t} を外せばよい。

\subsection{中心拡大の反復}

\begin{corollary}\label{cor:iterated-central}
$n\ge2$ とし、$C_0$ を定理~\ref{thm:central-ore} の $(1,1,2)$ 型代数とする。
連結局所有限次数付き代数の列 $C_0,C_1,\ldots,C_{n-2}$ があり、各 $j=1,\ldots,n-2$ について
\[
  t_j\in Z(C_j)_1\text{ は非零因子},\qquad
  C_j/(t_j)\cong C_{j-1}
\]
が成り立つとする。このとき $A=C_{n-2}$ は両側 noetherian な AS-regular代数で、
\[
 \gldim A=n+1,\qquad
 \ASind(A)=n+2,\qquad
 H_A(s)=\frac1{(1-s)^n(1-s^2)}
\]
を満たし、最小生成次数は $(1^n,2)$ である。従って定理~\ref{thm:ore} が適用され、
\[
 A=\langle A_1\rangle[w;\sigma,\delta],\qquad \deg w=2,
\]
と書ける。ここで $\langle A_1\rangle$ はquadratic Koszul $n$ 次元 AS-regular代数である。
\end{corollary}

\begin{proof}
$C_0$ の大域次元は $3$、AS indexは $4$、最小生成次数は $(1,1,2)$ であり、
定理~\ref{thm:central-ore} の証明で得た Ore 表示から
$H_{C_0}(s)=((1-s)^2(1-s^2))^{-1}$ である。
補題~\ref{lem:central-lift} を各段階に適用すると、
$j$ 段階目では大域次元は $3+j$、AS indexは $4+j$、
最小生成次数は $(1^{2+j},2)$ となり、Hilbert 級数には $(1-s)^{-j}$ が掛かる。
$j=n-2$ とすれば主張の数値と主定理のすべての仮定が得られる。
\end{proof}

\section{最小生成次数だけでは従わないこと}\label{sec:counterexample}

\begin{proposition}\label{prop:counterexample}
任意の \(n\ge1\) に対し、最小生成次数が \((1^n,2)\) である両側 noetherian
AS-regular domain \(A_n\) で、\(k\langle(A_n)_1\rangle\) 上の次数 \(2\) の
Ore 拡大ではないものが存在する。この例の大域次元は \(n+2\) である。
\end{proposition}

\begin{proof}
可換多項式環
\[
R=k[x_1,\ldots,x_n,u],\qquad |x_i|=1,\quad |u|=3
\]
上の次数 \(2\) の derivation \(D\) を
\[
D(x_1)=u,\qquad D(x_i)=0\ (i\ge2),\qquad D(u)=0
\]
で定め、\(A_n=R[w;D]\)、\(|w|=2\) とする。
従って
\[
wx_1-x_1w=u,\qquad wx_i=x_iw\ (i\ge2),\qquad wu=uw.
\]
Ore 拡大の自由加群表示により
\[
\{x_1^{a_1}\cdots x_n^{a_n}u^bw^c\mid a_i,b,c\ge0\}
\]
は \(k\)-基底をなす。\(u\) は regular central element であり、
\[
A_n/(u)\cong k[x_1,\ldots,x_n,w].
\]
右辺は大域次元 \(n+1\) の重み付き多項式環であるから、
補題~\ref{lem:central-lift} を \(e=3\) で適用すると、
\(A_n\) は両側 noetherian な大域次元 \(n+2\) の AS-regular 代数である。
Domain 性は、Ore 変数 \(w\) に関する最高次係数の積が非零であることから従う。

\(u=wx_1-x_1w\) なので \(u\) は最小生成元ではない。
一方
\[
(A_n)_1=\bigoplus_i kx_i,\qquad
(A_n)_2=k[x_1,\ldots,x_n]_2\oplus kw
\]
であるから、最小生成次数は正確に \((1^n,2)\) である。
\[
B_n:=k\langle(A_n)_1\rangle=k[x_1,\ldots,x_n],\qquad
H_{A_n}(t)=\frac1{(1-t)^n(1-t^2)(1-t^3)}.
\]
もし \(A_n=B_n[v;\sigma,\delta]\)、\(|v|=2\) なら
\(H_{A_n}=H_{B_n}/(1-t^2)\) となり、上式に反する。
\end{proof}

この例は、大域次元を指定しない場合の否定例である。
\(\gldim A=n+1\) まで仮定した一般の \((1^n,2)\) 型については、
\(\ASind A=n+2\) をさらに仮定すれば定理~\ref{thm:ore} が適用できる。
最小生成次数と大域次元だけからその AS index を導く一般命題は、
本稿の結果には含まれない。

\phantomsection
\addcontentsline{toc}{section}{参考文献}
\begin{thebibliography}{99}

\bibitem{LPWZ}
D.-M. Lu, J. H. Palmieri, Q.-S. Wu and J. J. Zhang,
\emph{Regular algebras of dimension 4 and their $A_\infty$-Ext-algebras},
Duke Math. J. \textbf{137} (2007), 537--584.
\href{https://arxiv.org/abs/math/0411497v1}{arXiv:math/0411497v1}.

\bibitem{CKS}
A. Chirvasitu, R. Kanda and S. P. Smith,
\emph{New Artin--Schelter regular and Calabi--Yau algebras via normal extensions},
Trans. Amer. Math. Soc. \textbf{372} (2019), 3947--3983.
\href{https://arxiv.org/abs/1706.05754v3}{arXiv:1706.05754v3}.

\bibitem{LWW}
L.-Y. Liu, S.-Q. Wang and Q.-S. Wu,
\emph{Twisted Calabi--Yau property of Ore extensions},
J. Noncommut. Geom. \textbf{8} (2014), 587--609.
\href{https://arxiv.org/abs/1205.0893v1}{arXiv:1205.0893v1}.

\bibitem{RZ}
D. Rogalski and J. J. Zhang,
\emph{Regular algebras of dimension 4 with 3 generators},
\href{https://arxiv.org/abs/1101.1998v2}{arXiv:1101.1998v2}.

\end{thebibliography}

\end{document}
